\section{Institutional Background}
\label{sec:institutional}

\subsection{Judicial orders as procurement shocks}
\label{subsec:litigation}

Brazil's constitutional right to health gives courts a central role
in disputes over access to medicines
\citep{ferraz2009right,wang2015right,cnj2019judicializacao}. In
S\~ao Paulo, patients can obtain court orders requiring the state
health administration to provide a specific medicine. The legal
process is initiated outside the procurement office: patients,
physicians, lawyers, public defenders, and judges determine whether
an order arrives, what medicine it covers, and when the agency must
respond. The institutional fact that matters for procurement is not
the legal doctrine itself, but the hard delivery obligation it
creates. The purchasing unit receives the mandate as an obligation to
source, not as a procurement option it selected.

This distinction is operational as well as legal. A procurement unit
can choose how to write a tender, when to aggregate ordinary demand,
and which procurement modality to use within statutory rules. It
cannot choose whether a plaintiff files a case, whether a judge grants
relief, or whether the order names a medicine that must be delivered
on a short schedule. Once the order arrives, the unit's choices are
downstream: how to comply, which suppliers to attract, and whether to
accept a costly outcome to avoid nondelivery.

This institutional sequence disciplines the paper's first comparison.
Court-mandated status is imposed on public buyers, not chosen to
obtain a favorable price, select a supplier, or improve a tender
outcome. The remaining concern is that the health-demand process
selects different patients or medicines into litigation. The design
addresses that concern with item, time, and buyer controls,
administrative-versus-litigated bounds, and mechanism evidence rather
than by claiming that selection into litigation is ignorable.

Judicial orders also change the procurement environment. They often
arrive with short deadlines and compliance measures, including daily
fines or fund seizure for nondelivery. The resulting pressure is
one-sided: the public buyer must deliver, while suppliers remain
voluntary participants who choose whether to bid and at what price.
Procurement notices may reference the court order or associated legal
process, so bidders can sometimes infer that the buyer is under
judicial urgency.

\subsection{Procurement on BEC}
\label{subsec:bec}

S\~ao Paulo's Department of Health buys medicines through BEC, the
state electronic procurement platform. BEC records item identifiers,
buyers, quantities, reference prices, negotiated prices, procurement
modalities, participating firms, winning suppliers, and tender
outcomes. Ordinary purchases, administrative urgent purchases, and
litigated urgent purchases are therefore observed within the same
procurement infrastructure.

Ordinary procurement can plan recurring demand, consolidate orders,
and wait for broader supplier participation. Urgent procurement
compresses this process. A buyer responding to a patient-specific
mandate may buy a smaller lot, repeat tenders more frequently, or
accept a thinner supplier set. Judicial urgency therefore has direct
sourcing consequences. Because BEC observes all three regimes on a
common platform, the data can separate same-firm pricing, quantities,
and winner switching.

\subsection{Administrative urgent demand}
\label{subsec:admin}

The administrative-request channel processes some patient demands
outside court. Patients can request medicines directly through the
health administration; accepted requests can generate urgent
individualized procurement without a judicial order. These purchases
share the operational pressure of individualized delivery but do not
carry the court-sanction regime described above.

The channel is valuable precisely because it is imperfect. It creates
urgent pharmaceutical demand outside court sanctions, but accepted
requests pass through scientific and cost-effectiveness screening and
can differ from litigated cases in item mix, feasibility, and scale.
The paper therefore uses administrative purchases as the closest
feasible urgent-procurement comparison, not as a clean
counterfactual. Ordinary purchases reflect planned demand;
administrative purchases reflect screened urgent requests without
court sanctions; litigated purchases reflect court-mandated urgent
delivery under possible fines, seizure, or liability.

This conservative use of the administrative channel is central. A
direct comparison would attribute the entire litigated-administrative
gap to sanctions, but the committee can admit cases that are more
feasible or more cost-effective. We instead ask whether a positive
gap remains after allowing for monotone selection into the
administrative channel, and through which procurement margin it
operates.
